\phantomsection
\addcontentsline{toc}{subsection}{Definiciones previas}
\subsection*{Definiciones previas}

\phantomsection
\addcontentsline{toc}{subsubsection}{Finanzas personales}
\subsubsection*{Finanzas personales}

La gestión de las finanzas personales se define como el proceso administrativo mediante el cual un individuo o núcleo familiar gestiona sus recursos económicos, abarcando el control de flujos de efectivo, ahorro y decisiones de inversión \parencite{UniversidadLibre01}. Este proceso requiere un enfoque metódico para mitigar desequilibrios presupuestales; en este contexto, la digitalización ha transformado la gestión financiera cotidiana al reducir las barreras de acceso y permitir la elaboración de presupuestos automatizados mediante herramientas tecnológicas \parencite{CESA01}.

De acuerdo con \textcite{Asobancaria01}, la adopción de servicios financieros digitales representa una oportunidad para mejorar la eficiencia en el seguimiento de gastos, lo cual impacta positivamente en la salud financiera del usuario. No obstante, la efectividad de estas herramientas está supeditada al nivel de alfabetización financiera de la población. Investigaciones recientes en el contexto colombiano indican que existe una brecha crítica en la comprensión de conceptos fundamentales como el interés compuesto y el valor del dinero en el tiempo, lo que limita la toma de decisiones informadas \parencite[p.~4]{Departamento01}.

Desde la perspectiva de la teoría de las finanzas personales, no basta con el acceso tecnológico; es imperativo el acompañamiento educativo para transformar la información en bienestar financiero \parencite{Inclusionfinan01}. Por lo tanto, se justifica el desarrollo de soluciones que no solo operen como gestores de datos, sino como sistemas de soporte a la decisión que ayuden a personas sin experiencia previa a alcanzar una seguridad económica a largo plazo \parencite{Departamento01}.

\phantomsection
\addcontentsline{toc}{subsubsection}{Nómadas digitales}
\subsubsection*{Nómadas digitales}

Los nómadas digitales constituyen un perfil emergente de trabajadores remotos que, gracias a la tecnología, desempeñan sus labores desde cualquier lugar del mundo. Como señalan \textcite{Recordar01} y refleja el \textcite{BBVA01}, este concepto define «un estilo de vida en el que, gracias a los avances tecnológicos, los profesionales pueden desempeñar su labor desde cualquier lugar y en cualquier momento». Se estima que actualmente existen unos 40 millones de nómadas digitales en el mundo, la mayoría adultos (30-39 años) con altos ingresos en el sector tecnológico.

En América Latina esta tendencia también crece: países como Colombia han implementado visas especiales para atraer a estos profesionales itinerantes. En este contexto, la población objetivo en Bogotá incluye a muchos migrantes y trabajadores móviles que carecen de una asesoría financiera tradicional. Para estos usuarios, gestionar ingresos variables y gastos internacionales presenta desafíos únicos. Por ello, el plan de negocio debe reconocer que los nómadas digitales valoran la flexibilidad financiera y requieren soluciones adaptables (por ejemplo, manejo de distintas monedas, ahorro para emergencias médicas o planes de retiro internacionales). Atender a este segmento implica ofrecer servicios de planificación financiera personalizados que contemplen su estilo de vida dinámico, ayudándolos a estabilizar sus finanzas pese a la movilidad geográfica.

\phantomsection
\addcontentsline{toc}{subsubsection}{Plan de negocio}
\subsubsection*{Plan de negocio}

Un plan de negocios es el documento rector donde se articula la propuesta de valor y la viabilidad de un emprendimiento. Constituye la ``tarjeta de presentación'' de la empresa ante inversores y colaboradores potenciales. En él se definen los objetivos estratégicos, la misión, el mercado meta, la estructura organizativa y los recursos necesarios, así como las proyecciones financieras a medio plazo. Según el \textcite{BBVA01}, un plan de negocios bien elaborado debe identificar la oportunidad de mercado, analizar su viabilidad técnica, económica y financiera, y describir las estrategias para convertir la idea en empresa real. 

En esencia, su objetivo principal es ayudar al emprendedor a reflexionar sobre todos los aspectos estratégicos de su proyecto. Esto incluye delinear un cronograma de resultados esperados y los pasos clave a seguir. Por ejemplo, se suelen estructurar apartados como resumen ejecutivo, análisis de mercado, descripción del servicio (planificación financiera), modelo de ingresos, plan operativo y proyecciones financieras (cuentas de resultados y flujo de caja). Cada sección contribuye a garantizar que el proyecto sea factible: desde cuantificar la demanda en Bogotá hasta prever inversiones en tecnología o capacitación. Un plan de negocios claro y detallado reduce la incertidumbre y provee un camino definido para tomar decisiones, convirtiendo una idea abstracta en una estrategia coherente para crear valor sostenible en el mercado.

\phantomsection
\addcontentsline{toc}{subsubsection}{Asesoramiento financiero}
\subsubsection*{Asesoramiento financiero}

El asesoramiento financiero es un servicio profesional cuyo objetivo es guiar a los individuos en la toma de decisiones económicas y patrimoniales, con el fin de alcanzar sus metas (compra de vivienda, jubilación, ahorro, etc.) y mejorar su salud financiera. En palabras del \textcite{BBVA02}, ``la función de un asesor es guiar a las personas para que puedan conseguir los objetivos que se hayan planteado, buscando siempre la mejora de la salud financiera de sus clientes''. Existen diferentes modalidades de asesoría según el vínculo del profesional:

\begin{itemize}
	\item \textbf{Asesor independiente:} No está ligado a ninguna entidad financiera. Este tipo de asesor actúa de forma imparcial, pudiendo recomendar productos de cualquier banco o fondo. Generalmente trabaja bajo contrato y no cobra comisiones por venta, lo que busca garantizar que el consejo sea objetivo.
	\item \textbf{Asesor dependiente o gestor comercial:} Trabaja dentro de una entidad financiera y suele ofrecer exclusivamente los productos de esa institución, percibiendo comisiones asociadas. Aunque su formación puede ser similar, su independencia está limitada por su relación con la empresa matriz.
\end{itemize}

Un buen asesor debe poseer conocimientos especializados en inversiones, gestión de riesgos y planificación patrimonial, y sobre todo debe adaptar el servicio al perfil de cada cliente. Entre las características señaladas por BBVA destaca la personalización del servicio: el asesor financiero analiza la capacidad de ahorro, nivel de endeudamiento, tolerancia al riesgo y situación socioeconómica del cliente, para diseñar recomendaciones a su medida. Asimismo, emplea un lenguaje claro y evita tecnicismos innecesarios, explicando abiertamente riesgos y beneficios de cada alternativa.

Todo ello busca generar confianza y transparencia con el cliente. Por otro lado, es importante mencionar que existe un marco regulatorio (p. ej. directivas MiFID en Europa) que protege al inversor minorista: estas normas exigen al asesor evaluar la conveniencia de los productos según el perfil del cliente, fomentando prácticas éticas y reduciendo conflictos de interés. En síntesis, el asesoramiento financiero combina experticio profesional con un enfoque centrado en el cliente, proporcionando orientación continua que construye seguridad y confianza en la gestión económica personal.

\phantomsection
\addcontentsline{toc}{subsubsection}{Herramientas tecnológicas para el asesoramiento financiero}
\subsubsection*{Herramientas tecnológicas para el asesoramiento financiero}

La revolución digital ha dado lugar a diversas herramientas tecnológicas que facilitan tanto la gestión cotidiana de las finanzas personales como el mismo proceso de asesoría. Entre las más relevantes se encuentran:

\begin{itemize}
	\item \textbf{Aplicaciones de gestión de gastos:} \textit{Apps} móviles como Toshl o Money Lover permiten registrar fácilmente ingresos y egresos diarios desde \textit{smartphones} y visualizar el presupuesto a través de gráficos interactivos. Estas aplicaciones animan al usuario a llevar un seguimiento disciplinado de sus gastos, al tiempo que hacen el proceso más ameno mediante interfaces lúdicas.
	\item \textbf{Plataformas agregadoras de cuentas bancarias:} Herramientas como BBVA Bconomy, Ahorro.net o Fintonic centralizan en un solo panel todas las cuentas y tarjetas del usuario. Al consolidar la información financiera, estas plataformas analizan patrones de gasto e ingresos en tiempo real y ofrecen recomendaciones personalizadas de ahorro. Por ejemplo, Bconomy evalúa la ``salud financiera'' del cliente y sugiere planes a medida, mientras que Ahorro.net se define como un ``asesor financiero de bolsillo'' que monitorea múltiples cuentas para aconsejar al usuario sobre cómo optimizar sus ahorros. Fintonic, por su parte, sincroniza automáticamente los movimientos bancarios y brinda alertas de movimientos y consejos personalizados de ahorro.
	
	\item \textbf{Chatbot asesores:} Se trata de plataformas automatizadas basadas en algoritmos que crean y gestionan carteras de inversión sin intervención humana continua. Según análisis de \textcite{BBVAResearch01}, estos servicios ofrecen asesoría financiera masiva a menor costo, orientada generalmente a inversión pasiva en fondos indexados (ETFs), pero adaptada al perfil de riesgo de cada cliente. Los \textit{robo-advisors} facilitan el acceso a la planificación de inversiones a segmentos con menores patrimios o con afinidad por canales digitales, al ofrecer transparencia en costos y recomendaciones objetivas.
	\item \textbf{Calculadoras y simuladores en línea:} Existen numerosas herramientas \textit{web} o móviles que ayudan a proyectar escenarios financieros básicos (simuladores de ahorros, préstamos, hipotecas, jubilación, etc.). Estas calculadoras apoyan la educación financiera práctica, al permitir evaluar rápidamente cómo variaciones en la tasa de interés, inflación o ingresos afectan el presupuesto personal. Si bien no sustituyen al asesor humano, complementan el proceso de planificación, empoderando al usuario para entender conceptos clave.
\end{itemize}

\phantomsection
\addcontentsline{toc}{subsubsection}{Alfabetización financiera}
\subsubsection*{Alfabetización financiera}

La alfabetización financiera se entiende, según \textcite{OECD01}, como el conjunto de conocimientos y entendimiento de conceptos financieros y sus riesgos, así como las habilidades y actitudes de aplicar dichos conocimientos y entender en orden a realizar decisiones efectivas a través de un rango de contextos financieros para mejorar el bienestar de los individuos y la sociedad. En este sentido, un individuo alfabetizado financieramente comprende conceptos básicos como presupuesto, interés, deuda y ahorro, y aplica esa comprensión para gestionar sus ingresos, gastos y ahorros. Estudios internacionales demuestran que, si se consideran causales los efectos de la alfabetización financiera sobre el comportamiento financiero, los costos de la ignorancia financiera son sustanciales \parencite[p.~17]{PMC01}. Por consecuente, aumentar la alfabetización financiera implica empoderar al usuario para enfrentar decisiones financieras cotidianas y extraordinarias con mayor seguridad y criterio.

\phantomsection
\addcontentsline{toc}{subsubsection}{Educación financiera}
\subsubsection*{Educación financiera}

La educación financiera es el proceso mediante el cual los individuos mejoran su comprensión sobre productos financieros, conceptos y riesgos, desarrollando habilidades y confianza para gestionar sus finanzas. Según la \textcite{OECD02}, este proceso combina información, instrucción y/o asesoría objetiva, de modo que las personas se hagan más conscientes de las oportunidades y riesgos financieros, aprendan dónde buscar ayuda y tomen acciones efectivas para mejorar su bienestar económico. En la práctica, la educación financiera básica incluye la enseñanza de presupuestos, manejo de deudas, evaluación de productos financieros y realización de inversiones, entre otros \parencite{Investopedia01}. Por ello, un servicio de planificación financiera para personas sin conocimiento previo debe incorporar componentes educativos claros y accesibles, que faciliten el aprendizaje continuo de conceptos financieros fundamentales.

\phantomsection
\addcontentsline{toc}{subsubsection}{Inclusión financiera}
\subsubsection*{Inclusión financiera}

La inclusión financiera se define como el proceso de promover el acceso asequible, oportuno y adecuado a una amplia gama de productos y servicios financieros regulados, ampliando su uso entre todos los segmentos de la sociedad. Dicho de otro modo, la inclusión financiera busca que individuos y negocios utilicen productos financieros (ahorro, crédito, pagos, seguros, etc.) que respondan a sus necesidades y les permitan gestionar riesgos y oportunidades \parencite[p.~22]{BancoDesarrollo01}. Este concepto cobra especial relevancia en poblaciones con baja alfabetización: la inclusión no solo abarca acceso físico a los bancos, sino también a herramientas y educación financiera que habiliten un uso efectivo de dichos servicios. Al integrar la inclusión financiera en el proyecto, se enfatiza la necesidad de reducir barreras (económicas, educativas y tecnológicas) para que más personas participen plenamente en la economía formal, mejorando así su estabilidad económica y posibilidades de crecimiento \parencite[p.~2]{OECD01}.

\phantomsection
\addcontentsline{toc}{subsubsection}{Tecnología financiera (\textit{FinTech})}
\subsubsection*{Tecnología financiera (\textit{FinTech})}

El término \textit{FinTech} (tecnología financiera) se refiere al uso de la tecnología y la innovación que tienen como propósito mejorar o automatizar los servicios financieros. En la práctica, esto incluye varios segmentos como los medios de pago y su infraestructura, proveedores de pagos, finanzas personales y elaboración de ahorros, entre otros \parencite{BancadeOportunidades02}. Las herramientas \textit{FinTech} facilitan el acceso a información financiera en tiempo real y permiten servicios personalizados a bajo costo, lo que es clave para usuarios con poco conocimiento previo. Por ejemplo, una \textit{app} de finanzas personales puede mostrar gráficos simples de gastos o enviar alertas de presupuesto, reduciendo la complejidad técnica. En el contexto del servicio propuesto, integrar tecnología financiera significa aprovechar plataformas digitales (\textit{apps}, simuladores, \textit{chatbots}, etc.) para apoyar la educación financiera y el asesoramiento personalizado.

\phantomsection
\addcontentsline{toc}{subsubsection}{Tecnología educativa (\textit{EdTech})}
\subsubsection*{Tecnología educativa (\textit{EdTech})}

El término \textit{EdTech} (tecnología educativa) engloba dos definiciones complementarias: la primera se refiere al estudio y solución de retos educativos mediante herramientas tecnológicas innovadoras; mientras que la segunda busca aprovechar el uso de plataformas interactivas para optimizar el proceso de enseñanza \parencite{UniversidadEuropea01}. Bajo el contexto de la baja educación financiera en nuestra nación, las \textit{EdTech} emergen como una solución prometedora para democratizar el acceso al conocimiento económico, facilitando la llegada de contenidos educativos a comunidades apartadas \parencite{LaSillaVacia01}, además de permitir una personalización del aprendizaje, logrando de esta forma optimizar procesos educativos volviendo los procesos educativos más eficientes \parencite{MundoPosgrado01}. Nuestro proyecto plantea una \textit{EdTech} con características de \textit{FinTech} que permita acceder a la educación financiera a poblaciones de menores estratos, jóvenes y nómadas digitales.

\phantomsection
\addcontentsline{toc}{subsubsection}{Bienestar financiero}
\subsubsection*{Bienestar financiero}

El bienestar financiero se describe generalmente como un estado derivado de la relación entre las variables de producción, empleo y distribución de la renta, maximizando el bienestar social mediante el crecimiento económico \parencite{Economipedia02}. Bajo este concepto, una persona puede cumplir completamente con sus obligaciones financieras presentes y futuras, sintiéndose segura respecto a su futuro económico y tomando decisiones que le permitan disfrutar de la vida. No se trata solo de tener ingresos, sino de tener control sobre el dinero día a día, capacidad para enfrentar imprevistos (fondo de emergencia) y libertad para alcanzar metas personales. Este concepto enfatiza tanto aspectos objetivos (cumplir pagos, niveles de ahorro o inversión) como subjetivos (confianza y satisfacción con la situación económica). Para el servicio de planificación, mejorar el bienestar financiero del usuario implica fortalecer su estabilidad económica doméstica mediante un presupuesto realista y opciones de inversión seguras.

\phantomsection
\addcontentsline{toc}{subsubsection}{Presupuesto personal}
\subsubsection*{Presupuesto personal}

El presupuesto personal es un plan financiero que detalla cómo se distribuirá el dinero a distintas categorías de gastos y ahorros durante un periodo determinado. Es una herramienta que permite visualizar claramente la relación entre ingresos, gastos y ahorro, facilitando la toma de decisiones económicas. Elaborar un presupuesto personal ayuda a identificar gastos innecesarios, asignar una parte de los ingresos a un plan de ahorros y planificar las prioridades en el uso del dinero \parencite{GrupoR01}. Para personas con bajo nivel financiero, aprender a realizar un presupuesto es fundamental, ya que les proporciona control sobre sus finanzas cotidianas y les enseña hábitos de ahorro y planificación práctica.

\phantomsection
\addcontentsline{toc}{subsubsection}{Aprendizaje Automático (\textit{Machine Learning})}
\subsubsection*{Aprendizaje Automático (\textit{Machine Learning})}

El \textit{Machine Learning} es una rama de la Inteligencia Artificial que permite a los sistemas aprender de los datos para realizar predicciones o tomar decisiones sin ser programados explícitamente para cada tarea. En el contexto de las \textit{FinTech}, se emplea para la segmentación de usuarios (\textit{Clustering}), la detección de patrones de gasto y la personalización de recomendaciones financieras. Para este proyecto, se explorará el uso de algoritmos de aprendizaje supervisado (como modelos de regresión para proyectar ahorros) y aprendizaje no supervisado (para agrupar perfiles de riesgo similares entre nómadas digitales y jóvenes). Esta capacidad tecnológica permitirá que la plataforma evolucione de un simple gestor de datos a un sistema inteligente de recomendación financiera.